\subsection{系统框架与软件算法演进}

视觉 SLAM 算法的演进是一个多种理论不断交叉融合的过程，从早期的滤波机制到现代的非线性优化，再到如今引入深度学习模型，其系统框架在精度、鲁棒性以及计算复杂度上不断寻求平衡。本节将从前端数据关联方式、后端状态估计框架以及深度学习的融合三个维度，对视觉 SLAM 的软件算法演进进行阐述。

\subsubsection{前端数据关联：间接法、直接法与半直接法}
视觉前端的主要任务是从图像流中提取特征并建立帧间数据关联，进而为后端提供初始的相机运动估计。根据数据关联方式的不同，前端主要分为间接法、直接法和半直接法三大类。

\textbf{间接法（特征点法）}是目前最为成熟和广泛应用的前端方案。该方法首先从图像中提取具有代表性的关键点（如 SIFT、SURF、ORB，以及 FAST 角点等），并计算相应的描述子；随后通过描述子匹配建立相邻帧之间的 2D-2D 或 2D-3D 数据关联，最后通过最小化重投影误差（Reprojection Error）求解相机位姿。间接法的优势在于对光照变化和图像模糊具有较强的鲁棒性，且提取的几何特征便于后端进行优化。典型的代表算法包括 PTAM 以其后风靡学术界与工业界的 ORB-SLAM 系列\cite{mur2015orb, mur2017orb, campos2021orb}。然而，特征提取和匹配过程极为耗时，往往构成了前端计算资源的瓶颈。

\textbf{直接法（Direct Method）}的提出旨在跳过耗时的特征提取环节。它基于灰度不变性假设（即空间中同一点在不同视角下的像素灰度保持不变），通过直接最小化光度误差（Photometric Error）来估计相机运动。直接法能够利用图像中的像素点（包括边缘和无明显角点区域），在弱纹理环境中表现出更好的鲁棒性。经典的直接法 SLAM 包括 DTAM\cite{newcombe2011dtam} 和 LSD-SLAM\cite{engel2014lsd}。尽管直接法省去了特征计算，但由于局部光度梯度的非凸性，其对相机的剧烈运动和光照突变极为敏感，容易陷入局部最优。

为了结合上述两者的优势，\textbf{半直接法（Semi-direct Method）}应运而生。半直接法通常仅提取 FAST 角点但不计算描述子，利用光流法或直接法追踪这些特征点，随后再结合特征对齐提升定位精度。它在一定程度上避免了直接法的非凸性问题，同时免去了特征点法冗长的描述子匹配开销。SVO\cite{forster2014svo} 是该领域的标杆算法，其凭借极高的前端运行帧率和良好的精度，成为微型飞行器（MAV）等计算受限平台的理想框架。

\subsubsection{后端状态估计：滤波法与优化法}
视觉 SLAM 的后端负责消除前端的累积漂移，通过融合多源测量（如 IMU、里程计等）维持全局一致的地图。按照状态估计类型的不同，主要分为滤波法和优化法。

\textbf{滤波法（Filter-based）}主要基于马尔科夫假设，利用贝叶斯滤波框架在当前时刻更新系统状态与协方差矩阵。早期的单目 MonoSLAM\cite{davison2007monoslam} 采用了扩展卡尔曼滤波（EKF）框架，将相机位姿与路标点共同纳入状态向量进行更新。但由于状态向量和协方差矩阵的大小随路标点的增加而膨胀，其计算量呈平方级增长（$O(N^2)$），严重限制了地图的大规模扩展。为解决此问题，Mourikis 等人提出了 MSCKF（Multi-State Constraint Kalman Filter）\cite{mourikis2007multi}，其核心思想是在状态向量中仅维护滑动窗口内的历史位姿，将针对同一路标的多帧观测投影到零空间消除路标状态，使计算复杂度仅与窗口大小相关。MSCKF 极大地提升了系统的实时性，成为了视觉惯性里程计（VIO）领域的经典方案，并广泛应用于 Google ARCore 等工业级产品。

随着图优化理论与稀疏矩阵线性代数求解器领域的成熟，\textbf{优化法（Optimization-based）}逐渐取代滤波法成为主流框架。优化法将所有或一段滑动窗口内的历史轨迹及路标点构建为一个超大的非线性最小二乘图，通过束调整（Bundle Adjustment, BA）全局解算位姿与路标。相比于 EKF-SLAM 在线性化点处的不一致问题，BA 允许进行多次重新线性化，并在全局范围内利用信息修正过去轨迹。当前主流系统如 ORB-SLAM 系列、OKVIS\cite{leutenegger2015keyframe} 以及 VINS-Mono\cite{qin2018vins} 均采用优化法。VINS-Mono 等将滑动窗口机制与预积分技术结合，通过边缘化（Marginalization）将早期状态信息转化为先验融入系统，在精度上表现优异。

\subsubsection{深度学习与 SLAM 的融合}
近年来，深度学习技术正在逐步渗透到 SLAM 的各个模块中。深层神经网络通过数据驱动的方式展现出了应对复杂动态场景、弱纹理区域以及剧烈光照变化的强大能力。
在前端提取层面，SuperPoint 等网络取代了传统的 FAST/ORB 提取器，提供更加稳健的特征点与描述子；在深层多视图几何理解方面，RAFT\cite{teed2020raft} 提出了基于全对偶场循环的架构用于准确预估光流。
到了系统层面，端到端的深度 SLAM（例如 DROID-SLAM\cite{zhu2020droid}）利用循环神经网络（RNN/GRU）迭代更新相机位姿与深度图，能够在大运动与严重模糊条件下鲁棒地完成跟踪建图。尽管深度学习 SLAM 展现出了高鲁棒性与语义感知潜力，但其依赖高性能 GPU 运算的特性，极大地限制了其在功耗和算力受限边缘芯片平台上的大规模部署。

\subsubsection{后端巨大计算开销与硬件加速的必然性}
综上所述，虽然基于优化法（BA）的 SLAM 在精度与鲁棒性方面表现优异，但在嵌入式系统中它却带来了沉重的计算负担。一次典型的后端 BA 优化需要通过计算繁琐的雅可比求导构建巨大的海森矩阵（Hessian Matrix），并利用 Schur 补等机制求解线性方程组。这一复杂的矩阵迭代过程在移动终端（如基于 ARM 处理器的无人机或 AR 设备上）产生了几十到上百毫秒量级的延迟，成为制约系统实时性与能效比的瓶颈。

针对 SLAM 的高昂算力需求，学术界首先在前端探索了专用领域的硬件加速方案。如 Liu 等人设计的 eSLAM\cite{liu2019eslam} 与 Fang 等人提出的前端加速架构\cite{fang2019fpga}，皆通过 FPGA 全流水线成功实现了高速特征检测与图像特征匹配，大幅降低了系统的时延与功耗。然而，当前端的处理能力得到了成倍的提升后，原本就吃紧的后端优化阶段面临着更加密集的特征点输入约束。这意味着必须突破通用处理器的计算瓶颈，针对基于 BA 的稀疏矩阵运算法则设计更深层级与更大规模数据的专用处理架构，方能在低功耗平台满足 SLAM 复杂算法迭代的需要。因此，将目光聚焦于通过软硬协同或者底层算子加速来提升 SLAM 后端运算效率，已成为了突破嵌入式领域自动导航系统性能天花板的必由之路。
